% !TITLE 过邺下吊高神武
% !AUTHOR 袁枚
% !DYNASTY 明清
% !GENRE 七言律诗
% !ORDER 9

\begin{poem}
唱罢阴山敕勒歌\textsuperscript{1}，英雄涕泪老来多。
生持魏武朝天笏\textsuperscript{2}，死授条侯\textsuperscript{3}杀贼戈。
六镇\textsuperscript{4}华夷\textsuperscript{5}传露布\textsuperscript{6}，九龙\textsuperscript{7}风雨聚漳河\textsuperscript{8}。
祇今\textsuperscript{9}尚有清流月，曾照高王\textsuperscript{10}万马过。
\end{poem}

\begin{annotations}
\notetext{1}{阴山敕勒歌：即我们在南北朝部分读过的《敕勒歌》。北齐神武帝高欢病重时，令大将斛律金唱此歌，高欢和之，涕泪纵横。}
\notetext{2}{朝天笏：朝见天子时手持的笏板。笏（hù），古代大臣上朝时拿的狭长板子，用于记事。这里代指高欢掌控朝政的权臣地位。}
\notetext{3}{条侯：指西汉条侯周亚夫，汉文帝、景帝时名将，以治军严整、平定七国之乱著称。将高欢比作周亚夫。}
\notetext{4}{六镇：北魏在北方边境设立的六个军事重镇（沃野、怀朔、武川、抚冥、柔玄、怀荒）。高欢出身怀朔镇，从边镇士卒起家。}
\notetext{5}{华夷：汉族（华）和各少数民族（夷）。泛指各民族。}
\notetext{6}{露布：古代一种不封口的告捷文书，捷报一路传抄、公开张贴以广而告之。}
\notetext{7}{九龙：指邺城的九龙殿。邺城是高欢建立的北齐政权的都城。}
\notetext{8}{漳河：河流名，流经邺城附近。}
\notetext{9}{祇今：至今。“祇”此处读作qí。}
\notetext{10}{高王：指高欢，北齐建立后被追尊为神武帝，亦称“高王”。}
\end{annotations}

\begin{appreciation}
这是清代诗人袁枚路过邺城时凭吊北齐神武帝高欢的一首七律，也是把我们在南北朝读过的那首《敕勒歌》重新唱响的诗之一。

诗从《敕勒歌》写起："唱罢阴山敕勒歌，英雄涕泪老来多。"《敕勒歌》我们读过：天苍苍，野茫茫，风吹草低见牛羊。高欢病重时，让大将斛律金唱这首歌，自己跟着和，泪流满面。一个打了一辈子仗的枭雄，临终前被这首草原上的歌击溃。袁枚用这一笔，就把英雄写成了人——英雄也会老，老了也会哭。

"生持魏武朝天笏，死授条侯杀贼戈。"魏武指曹操——也是当了一辈子权臣、到死没有称帝的人；条侯是汉代的周亚夫，以军纪严明出名。生与死对举：生前握着朝笏执掌朝政，死后被人比作名将。两句写尽高欢的一生。

"六镇华夷传露布，九龙风雨聚漳河。"六镇是北魏北境的六个军镇，高欢就从怀朔镇的小兵起家；露布是报捷的文书；九龙殿和漳河都在邺城，高欢定都的地方。起兵时捷报飞传，各族军民追随，最后万里江山都在脚下——一生的风云际会，压进两句诗里。

末联最有味道："祇今尚有清流月，曾照高王万马过。"至今只剩下那轮清清的月亮，曾照着高王的万马千军过漳河。人没了，马没了，城也荒了，只有月亮还在。李白《把酒问月》写过："今人不见古时月，今月曾经照古人。"一样的月亮，一样的感慨。中国诗写物是人非，最常用的就是这一招，却总能用得让人心里一沉。

袁枚是性灵派，主张诗写真性情。他凭吊高欢，不写权谋不写功业，专挑他听歌落泪的那一刻写——英雄的软弱处，才是诗人下笔的地方。邺城也是曹操的城，我们读过《观沧海》，曹操也是马上打天下的人。两个朝代的英雄在同一条漳河边点兵，到头来都只剩一轮月亮。

你背历史的时候，那些大人物的名字背了一遍又一遍——其实他们都是过漳河的风。所以我们对你不抱什么大指望，不指望你当英雄。英雄老了也会流泪，你把眼前的日子过好，比什么都强。
\end{appreciation}
